% Methods text supporting the production Figure 4.

\subsection{Retinal-motion replay and spatial information analysis}

We sampled 250-ms fixation histories from continuous DDPI recordings. Eye
position was placed on a uniform raw-sampling grid, filtered with the
zero-phase analysis IIR selected by the eye-position power-spectrum audit
(20-Hz passband, 30-Hz stopband), and sampled at the model's native 240-Hz
input rate. Each trajectory translated a natural image through the same
renderer used for model evaluation. The matched stabilized condition held gaze
at the frame corresponding to the model's rounded mean first-layer peak-energy
lag (15 frames; 62.5 ms before prediction). Measured and stabilized histories
therefore shared the retinal frame with greatest average first-layer weight
and differed in preceding motion. The 42-dimensional behavioral input was held
at zero in both conditions.

For each output unit, we evaluated its fitted spatial readout at every
available position to obtain a nonnegative activation map $r(x)$. Single-spike
information was
\[
  I_{\mathrm{SSI}} = \sum_x \frac{r(x)}{\sum_{x'}r(x')}
  \log_2\!\left(\frac{r(x)}{\bar r}\right),
\]
where $\bar r$ is the spatial mean. SSI is invariant to uniform multiplicative
gain. Panel B used all 725 checkpoint-available exact readouts, 40 natural
images, and 200 filtered fixation histories without microsaccade
stratification. Path length was measured after filtering over the same fixed
window. Boxes report distributions over exact units; population intervals were
obtained by crossed resampling of images, traces, and units.

\subsection{Exact-unit spatiotemporal tuning and retinal-power engagement}

We measured model SF$\times$TF tuning with controlled drifting gratings at the
native 240-Hz rate through every checkpoint-native
$(\mathrm{session},\mathrm{cid})$ readout. The assay contained nine spatial
frequencies (1.07--15.94 cycles/deg), 14 nonzero temporal frequencies
(1--90.51 Hz) plus a static condition, 18 motion directions, and 24 endpoint
phases at 0.25 contrast; it was repeated at 0.125 contrast. The primary
response was phase-averaged expected count per 1/240-s bin minus an explicit
gray blank. Yu R0 and R1 surfaces were fit to every dynamic SF$\times$TF cell
at the preferred motion direction. The strict exemplar audit required exact
identity, a complete grid, 24-versus-12-phase convergence, positive drive, an
interior coherent raw peak, full-grid fit $R^2\geq0.80$, raw/fit peak
agreement, contrast repeatability, reliable non-boundary recorded-neuron SF,
and recorded/model SF agreement. Recorded gratings constrained biological SF
but did not vary TF; all TF values are therefore model measurements.

Panel D used two units passing every strict gate. Panels E--H used a separate,
explicit all-unit contract containing every checkpoint-available exact readout
with a finite Yu fit (725 units). This contract preserves one-to-one
$(\mathrm{session},\mathrm{cid})$ identity and retains, labels, and reports the
units that fail stricter validation rather than treating them as validated.

For Panels C and F, we used the Kuang/Rucci translation identity to combine
radial natural-image power with the phase spectrum of filtered real fixations.
Drift-rich and rapid-transient quartiles were defined within animal from each
trace's nonzero-TF spectral centroid without reading neural or model responses.
Each trace's TF$>0$ spectrum was normalized to unit mass before
animal-balanced averaging; the two conditions therefore compare spectral
shape. The exact TF$=0$ carrier and TF$>0$ dynamic fractions were independently
verified to sum to one for every spatial Fourier mode.

For Panel G, spectra were recomputed from the directly rendered finite movies
for the same 40 images and 200 traces in the response matrix. Joint passband
power was the projection of each movie's
SF$\times$TF$\times$orientation power onto each unit's tuning weights.
Fixation histories were ranked within unit by measured-minus-stabilized
passband power. Rate and SSI modulation were calculated from the matched model
responses, not from the spectral predictor. The comparison with path length
used paired within-unit Spearman correlations and a unit bootstrap of their
difference.

\subsection{Cumulative trained-readout analysis}

For Panel H, we selected unit--movie pairs in the top 20\% of each unit's
passband-power distribution and replayed the matched measured and stabilized
movies through the unmodified model. Trained output logits were accumulated
over the S1-plus-phase, S2, and S3 feature groups, with the ordinary softplus
link applied after each cumulative sum. The final cumulative map was required
to match the ordinary model output numerically. Temporal modulation was
calculated after removing each location's temporal mean and normalized by the
movie-wide mean response; spatial sharpening was SSI of the mean-normalized
activation map. Thus neither displayed quantity can be generated by a uniform
output gain. Boxes show unit distributions and intervals resample exact units.
The production release requires at least 100 crossed image--fixation movies;
smaller runs are labeled smoke tests.
